\documentclass[11pt]{article}

\usepackage[margin=1in]{geometry}
\usepackage{amsmath,amssymb,amsthm,mathtools}
\usepackage{graphicx}
\usepackage{xcolor}
\usepackage{microtype}
\usepackage{enumitem}
\usepackage{booktabs}
\usepackage[hidelinks]{hyperref}

\newcommand{\C}{\mathbb C}
\newcommand{\R}{\mathbb R}
\newcommand{\conv}{\operatorname{conv}}
\newcommand{\diag}{\operatorname{diag}}
\newcommand{\spec}{\sigma}
\newcommand{\cB}{\mathcal B}
\newcommand{\cC}{\mathcal C}
\newcommand{\cE}{\mathcal E}
\newcommand{\cS}{\mathcal S}
\newcommand{\cM}{\mathcal M}
\numberwithin{equation}{section}

\newtheorem{theorem}{Theorem}
\newtheorem{corollary}[theorem]{Corollary}
\newtheorem{lemma}[theorem]{Lemma}
\theoremstyle{definition}
\newtheorem{remark}[theorem]{Remark}

\title{An exact algebraic envelope for four-by-four Ritz sets\\
with an interior eigenvalue}
\author{Candidate partial result for arXiv:2004.05288}
\date{17 August 2026}

\begin{document}
\maketitle

\begin{center}
\fbox{\parbox{0.92\textwidth}{
\textbf{Status.} Candidate substantial partial result, likely valid, pending
human review.  The theorem completely treats the first missing $4\times4$
case under the stated nondegeneracy assumptions; the source's broad questions
for arbitrary dimension and for $k\geq3$ remain open.}}
\end{center}

\begin{abstract}
Dela Rosa and Woerdeman ask for a characterization of
$\conv\cB_A(\mu)$ when a normal matrix has eigenvalues in the interior of its
numerical range.  We derive an explicit real-rooted support polynomial in all
dimensions.  When $A$ has four distinct noncollinear eigenvalues and
$\mu\notin\spec(A)$, its weight polytope is a segment and the polynomial has
bidegree at most $(2,2)$.  In every direction the exact support is attained at
an endpoint root, an interior simple stationary root, or an interior multiple
root.  These candidates are algebraic; generically the two interior searches
reduce to quartics.  Intersecting the resulting supporting half-planes gives
$\conv\cB_A(\mu)$ exactly.  This covers spectra with an interior eigenvalue,
where the polygonal formula of the source can fail.
\end{abstract}

\section{Source question and scope}

The source is K.~L. Dela Rosa and H.~J. Woerdeman, \emph{Location of Ritz
values in the numerical range of normal matrices}, arXiv:2004.05288,
published in \emph{Linear and Multilinear Algebra} 69 (2021), 2749--2778
\cite{DRW-location}.  Section~7, page~30 of the arXiv PDF, asks:

\begin{quote}
Develop similar results for $k$-Ritz sets, $3\leq k<n-1$.  Also consider a
normal matrix with some eigenvalues in the interior of its numerical range;
in that case it is less clear how to characterize $\conv\cB_A(\mu_1)$.
\end{quote}

The source's Figure~2 shows that its earlier polygonal formula can fail as
soon as one eigenvalue is interior.  The following is the required source
screenshot; it includes both open problems and that numerical example.

\begin{figure}[p]
  \centering
  \includegraphics[width=0.98\textwidth]{figures/open_problem_crop.png}
  \caption{The open problems and Figure~2 in \cite[p.~30]{DRW-location},
  rendered from the stored source PDF.}
  \label{fig:source}
\end{figure}

This packet solves the complete $4\times4$ support-envelope problem when the
eigenvalues are distinct and noncollinear and the prescribed first Ritz value
is not itself an eigenvalue.  It also gives a structural formula in arbitrary
dimension.  It does not solve the higher-$k$ question, the singular case
$\mu\in\spec(A)$, or the search for a comparably simple geometry in dimensions
$n\geq5$.

\clearpage
\section{Setting}

For $A\in\C^{n\times n}$ and $\mu\in W(A)$, define
\[
 \cB_A(\mu)=\left\{\nu\in\C:
 V^*AV=\begin{bmatrix}\mu&*\\0&\nu\end{bmatrix}
 \text{ for some }V\in\C^{n\times2},\ V^*V=I_2\right\}.
\]
Unitary similarity lets us take
$A=\Lambda=\diag(\lambda_1,\ldots,\lambda_n)$.  Put
\[
 \cC_\Lambda(\mu)=\left\{t\in\R_+^n:
 \sum_i t_i=1,\quad \sum_i t_i\lambda_i=\mu\right\}.
\]
For $t\in\cC_\Lambda(\mu)$ set $u_i=\sqrt{t_i}$.  The condition that
$[u,v]$ give an upper triangular compression is precisely
\[
 v\perp u,\qquad v\perp\Lambda u.
\]
Consequently, if the columns of $F_t$ form an orthonormal basis of
$\{u,\Lambda u\}^{\perp}$, the standard parametrization in the source is
\begin{equation}
 \cB_A(\mu)=\bigcup_{t\in\cC_\Lambda(\mu)}
 W(F_t^*\Lambda F_t).
 \label{eq:param}
\end{equation}
The choice $u_i\geq0$ loses no generality: phases can be removed by a diagonal
unitary commuting with $\Lambda$.

For a direction $\theta\in\R$, write
\[
 h_i(\theta)=\Re(e^{-i\theta}\lambda_i),\qquad
 H_\theta=\diag(h_1(\theta),\ldots,h_n(\theta)).
\]
The support of $W(F_t^*\Lambda F_t)$ in that direction is the largest
eigenvalue of $F_t^*H_\theta F_t$.

\section{Idea of the proof}

The apparent obstruction is that \eqref{eq:param} is a continuum of numerical ranges.
The key observation is that every directional Hermitian compression has a
single explicit characteristic polynomial.  Cauchy--Binet expands the
determinant using $(n-2)$-row minors of $F_t$.  Because $F_t$ is completed by
the two-dimensional space spanned by $u$ and $\Lambda u$, complementary
minors turn each coefficient into the elementary weight
\[
 t_i t_j|\lambda_i-\lambda_j|^2.
\]
Thus there is no hidden choice of basis $F_t$.

For four eigenvalues the three real barycentric constraints leave only one
weight parameter $s$.  The support is the maximum of the larger root of a
quadratic $P_\theta(s,x)$.  Ordinary one-variable calculus now exhausts every
maximizer: an endpoint, a differentiable stationary root, or a collision of
the two roots.  Those three cases are expressed by two endpoint quadratics and
the systems $P=P_s=0$ and $P=P_x=0$.  The proof below makes both reductions
exact.  The remaining broad problem is not existence of an algorithm, but a
simple higher-dimensional geometric description.

\section{The support polynomial in all dimensions}

Define
\begin{align}
 \delta(t)&=\sum_{1\leq i<j\leq n}
 t_it_j|\lambda_i-\lambda_j|^2,\label{eq:delta}\\
 P_\theta(t,x)&=\sum_{1\leq i<j\leq n}
 t_it_j|\lambda_i-\lambda_j|^2
 \prod_{\ell\notin\{i,j\}}(x-h_\ell(\theta)).\label{eq:poly}
\end{align}

\begin{theorem}[Directional compression polynomial]\label{thm:polynomial}
Let $\Lambda=\diag(\lambda_1,\ldots,\lambda_n)$ and let
$\mu\in W(\Lambda)\setminus\spec(\Lambda)$.  For every
$t\in\cC_\Lambda(\mu)$,
\begin{equation}
 \det\bigl(xI_{n-2}-F_t^*H_\theta F_t\bigr)
   =\frac{P_\theta(t,x)}{\delta(t)}.
 \label{eq:char}
\end{equation}
In particular all zeros of $P_\theta(t,\cdot)$ are real, and
\begin{equation}
 h_{\conv\cB_A(\mu)}(\theta)
 =\max_{t\in\cC_\Lambda(\mu)}
   \max\{x\in\R:P_\theta(t,x)=0\}.
 \label{eq:support-general}
\end{equation}
\end{theorem}

\begin{proof}
Let $Q=[u,\Lambda u]$.  Since $\mu\notin\spec(\Lambda)$, the two columns of
$Q$ are linearly independent for every feasible $t$: dependence would force
all positive-weight eigenvalues to equal $\mu$.  Hence
\[
 E=Q(Q^*Q)^{-1/2}
\]
has orthonormal columns and $U=[F_t,E]$ is unitary.  Since
$xI-F_t^*H_\theta F_t=F_t^*(xI-H_\theta)F_t$, Cauchy--Binet gives
\begin{equation}
 \det(xI-F_t^*H_\theta F_t)
 =\sum_{|I|=n-2}|\det F_t[I,:]|^2
   \prod_{\ell\in I}(x-h_\ell(\theta)).
 \label{eq:cb}
\end{equation}
Complementary minors of the unitary $U$ have equal modulus.  If
$I^c=\{i,j\}$, then
\begin{align*}
 |\det F_t[I,:]|^2
 &=|\det E[\{i,j\},:]|^2\\
 &=\frac{|\det Q[\{i,j\},:]|^2}{\det(Q^*Q)}
 =\frac{t_it_j|\lambda_i-\lambda_j|^2}{\det(Q^*Q)}.
\end{align*}
Moreover,
\begin{align}
 \det(Q^*Q)
 &=\sum_i t_i|\lambda_i|^2-\left|\sum_i t_i\lambda_i\right|^2\notag\\
 &=\sum_{i<j}t_it_j|\lambda_i-\lambda_j|^2
 =\delta(t)>0.\label{eq:gram}
\end{align}
Substitution in \eqref{eq:cb} proves \eqref{eq:char}.  Real-rootedness follows because the
left-hand side is the characteristic polynomial of a Hermitian matrix.
Finally, \eqref{eq:param} and the variational characterization of the largest Hermitian
eigenvalue give \eqref{eq:support-general}.  Taking the convex hull does not change a support
function.
\end{proof}

\section{Exact four-by-four envelope}

Assume now that $n=4$ and the four eigenvalues are distinct and not all
collinear.  The real $3\times4$ matrix with rows
\[
 (1,1,1,1),\qquad (\Re\lambda_i)_{i=1}^4,
 \qquad(\Im\lambda_i)_{i=1}^4
\]
has rank three.  Thus every nonempty $\cC_\Lambda(\mu)$ is a possibly
degenerate segment
\[
 t(s)=t^0+s d,\qquad a\leq s\leq b.
\]
For fixed $\theta$ abbreviate
\[
 p(s,x)=P_\theta(t(s),x),
 \qquad r_+(s)=\max\{x\in\R:p(s,x)=0\}.
\]
The polynomial $p$ has degree at most two in each variable.  Its leading
coefficient as a polynomial in $x$ is $\delta(t(s))>0$.

Define three sets of real candidate values:
\begin{align*}
 \cE_\theta&=\{x:p(a,x)=0\text{ or }p(b,x)=0\},\\
 \cS_\theta&=\{x:\exists s\in[a,b],\ p(s,x)=p_s(s,x)=0\},\\
 \cM_\theta&=\{x:\exists s\in[a,b],\ p(s,x)=p_x(s,x)=0\}.
\end{align*}
If $a=b$, interpret $\cE_\theta$ as the roots at that single point.

\begin{theorem}[Exact $4\times4$ interior-spectrum envelope]
\label{thm:four}
Let $A$ be a normal $4\times4$ matrix with four distinct noncollinear
eigenvalues, and let $\mu\in W(A)\setminus\spec(A)$.  Then, in every direction
$\theta$,
\begin{equation}
 h_{\conv\cB_A(\mu)}(\theta)
 =\max(\cE_\theta\cup\cS_\theta\cup\cM_\theta).
 \label{eq:support-four}
\end{equation}
Consequently
\begin{equation}
 \conv\cB_A(\mu)=
 \bigcap_{\theta\in\R}
 \left\{z\in\C:\Re(e^{-i\theta}z)
 \leq h_{\conv\cB_A(\mu)}(\theta)\right\},
 \label{eq:halfplanes}
\end{equation}
where each support value is determined by polynomial systems of bidegree at
most $(2,2)$.
\end{theorem}

\begin{proof}
If the segment is a point, \eqref{eq:support-four} is immediate.  Suppose $a<b$.  The two roots
of $p(s,\cdot)$ are eigenvalues of the continuous family of $2\times2$
Hermitian compressions described by Theorem~\ref{thm:polynomial}; equivalently,
the quadratic formula shows directly that $r_+$ is continuous.  It therefore
has a maximizer $s_*$ on $[a,b]$.

If $s_*$ is an endpoint, the value $r_+(s_*)$ lies in $\cE_\theta$.  Suppose
$s_*\in(a,b)$.  If $r_+(s_*)$ is a simple root, the implicit-function theorem
applied to $p(s,r_+(s))=0$ gives
\[
 r_+'(s_*)=-\frac{p_s(s_*,r_+(s_*))}
                  {p_x(s_*,r_+(s_*))}=0,
\]
so the value lies in $\cS_\theta$.  If it is not simple, it lies in
$\cM_\theta$.  Hence the right side of \eqref{eq:support-four} is at least the
support in \eqref{eq:support-general}.

Conversely, every $x$ in any of the three candidate sets is a root of
$p(s,\cdot)$ for some feasible $s$.  It is therefore at most $r_+(s)$, which
is at most the support.  This proves \eqref{eq:support-four}.  Formula
\eqref{eq:halfplanes} is the standard
supporting-half-plane reconstruction of a compact convex subset of $\C$.
\end{proof}

\begin{corollary}[Generic quartic computation]\label{cor:quartic}
Write
\[
 p(s,x)=A(x)s^2+B(x)s+C(x)
       =a(s)x^2+b(s)x+c(s).
\]
Generically, the simple stationary candidates in Theorem~\ref{thm:four} are
obtained from
\begin{equation}
 B(x)^2-4A(x)C(x)=0,\qquad
 s=-\frac{B(x)}{2A(x)}\in(a,b),
 \label{eq:quartic-stat}
\end{equation}
and the multiple-root candidates from
\begin{equation}
 b(s)^2-4a(s)c(s)=0,\qquad
 x=-\frac{b(s)}{2a(s)},\quad s\in(a,b).
 \label{eq:quartic-mult}
\end{equation}
Both displayed polynomial equations have degree at most four.  Together with
the roots of the endpoint quadratics and feasibility filtering, they give the
exact support.  In degenerate cases one uses the original systems
$p=p_s=0$ and $p=p_x=0$; if a stationary component persists, its simple root
is constant on that component, and its boundary value is already captured by
an endpoint or a multiple-root candidate.
\end{corollary}

\begin{proof}
At a simple stationary solution with $A(x)\neq0$, the equation
$p_s=2A(x)s+B(x)=0$ gives the stated $s$, and substitution into $p=0$ gives
$B^2-4AC=0$.  Each of $A,B,C$ has degree at most two in $x$, so this equation
has degree at most four.  A multiple root in $x$ is exactly a zero of the
quadratic discriminant $b^2-4ac$, also of degree at most four in $s$; here
$a(s)=\delta(t(s))>0$.  Vanishing leading coefficients are handled without
division by retaining the two original polynomial equations.  On any
positive-dimensional simple-root component of $p=p_s=0$, implicit
differentiation makes the root constant.  The final assertion follows by
continuity and the classification in Theorem~\ref{thm:four}.
\end{proof}

\section{Verification and failed simplifications}

The included script \texttt{code/verify\_envelope.py} constructs the weight
segment, evaluates \eqref{eq:poly}, and compares its two roots with a direct orthogonal
nullspace compression.  With deterministic seed 20260817 it passed 600 such
comparisons; the worst absolute root discrepancy was
$1.443\times10^{-15}$.  It also numerically optimized directional supports
and checked $p_s=0$ at 45 interior maximizers, with worst normalized residual
$5.434\times10^{-10}$.

For the source-like example
\[
 (\lambda_1,\lambda_2,\lambda_3,\lambda_4)
 =(i,0,1,0.35+0.33i),\qquad \mu=0.10+0.30i,
\]
the exact segment envelope exceeds the two endpoint ellipses by
$0.055372249026$ in one tested direction, and 66 of 361 tested directions
have an interior maximizing weight.  Thus endpoint triangles alone cannot
recover the Figure~2 phenomenon.  A different tempting relaxation---slicing
the unrestricted compression-polynomial coefficient polytope by the fixed
root---overestimates support by as much as $0.196837808229$ in this example.
The code is a consistency audit, not part of the proof.

\section{Novelty check, limitations, and review recommendation}

A bounded search through 17 August 2026 used the exact source title and
authors, forward-citation records, and the phrases ``interior eigenvalues'',
``$k$-Ritz sets'', and ``$\cB_A(\mu)$''.  The closest later work found is Dela
Rosa and Woerdeman's continuity paper \cite{DRW-continuity}, which proves
continuity of associated submatrix and Ritz-set maps, not the convex-envelope
formula here.  Dela Rosa's 2021 dissertation repeats the open directions.
Searches for algebraic Ritz-tuple characterizations found Cauchy-matrix
criteria, but no fixed-$\mu$ support-polynomial or quartic envelope theorem.
No matching solution was found; this is a bounded negative search, not a proof
of novelty.

The exact remaining limitations are:
\begin{enumerate}[leftmargin=2em]
\item Theorem~\ref{thm:four} excludes $\mu\in\spec(A)$, where $Q$ may lose
rank and a separate analysis is required.
\item For $n\geq5$, Theorem~\ref{thm:polynomial} gives an exact optimization
over a higher-dimensional polytope.  Facewise Lagrange-multiplier equations
give a finite semialgebraic procedure, but not a concise geometric answer.
\item The source's separate question about simultaneous $k$-Ritz sets for
$3\leq k<n-1$ remains open.
\end{enumerate}

The verdict is \textbf{candidate partial, likely valid}.  Human review should
focus on the minor identity in \eqref{eq:cb}--\eqref{eq:gram} and the exhaustive maximum
classification in Theorem~\ref{thm:four}.  The degenerate-component sentence
in Corollary~\ref{cor:quartic} is a third review point.  The packet deliberately
does not promote the result as a full answer to either broad source question.

\begin{thebibliography}{9}
\bibitem{DRW-location}
K.~L. Dela Rosa and H.~J. Woerdeman,
\emph{Location of Ritz values in the numerical range of normal matrices},
Linear and Multilinear Algebra \textbf{69} (2021), 2749--2778.
doi:10.1080/03081087.2020.1761280; arXiv:2004.05288.

\bibitem{DRW-continuity}
K.~L. Dela Rosa and H.~J. Woerdeman,
\emph{Continuity of submatrices and Ritz sets associated to a point in the
numerical range}, Linear Algebra Appl. \textbf{624} (2021), 1--13.
doi:10.1016/j.laa.2021.03.031.
\end{thebibliography}

\end{document}
